% !TeX root = ../教材题库_全部条目_转写工作稿.tex
% Source: XJ-B1-C05-HW-P202-Y0419
% 第5章 / 5.5 三角函数模型的简单应用 / 教材 p.202 / PDF 第209页
% OCR未校对：公式与文字均可能有误；未完成数学TeX转写；禁止直接发给学生
% 原题及相邻原解截图（核对用）：../题目原图/XJ-B1-C05-HW-P202-Y0419-01.webp
% 本文件只在转写完成后作为题干使用；OCR 条目尚不是可靠题干。
\begin{bankitem}{习题5.5 第3题}{第5章 / 5.5 三角函数模型的简单应用 / 教材 p.202 / PDF 第209页}{XJ-B1-C05-HW-P202-Y0419}{OCR未校对：公式与文字均可能有误；未完成数学TeX转写；禁止直接发给学生}
\begin{ocrdraft}
[教材p.202 / PDF第209页]
\par
3．如图为一个公路隧道，隧道口截面为正弦型
\par
曲线，已知隧道跨径为\&4m,最高点离地面4．5m．
\par
《4巧m
\par
（1）若设正弦型曲线与路面的左交点为原点0，
\par
试求出该曲线的函数解析式；
\par
》一42m—l
\par
（2）如果路面宽度为4．2m，试求出公路边缘距
\par
隧道顶端的高度．
\par
（第3题）
\end{ocrdraft}
\end{bankitem}
